\documentclass[10pt,a4paper]{amsart}
\usepackage[margin=19mm]{geometry}
\usepackage{amsmath,amssymb,mathtools}
\usepackage[scr=rsfs,cal=euler]{mathalfa}
\usepackage{xfrac}
\usepackage[unicode,hidelinks,bookmarks=false]{hyperref}
\newcommand{\ie}{i.e.\ }
\newcommand{\cf}{cf.\ }
\newcommand{\resp}{respectively\ }
\newcommand{\Subsection}{\subsection}
\providecommand{\nobreakdash}{\nobreak\hbox{-}}
\setlength{\emergencystretch}{2em}
\multlinegap0pt
\usepackage{amssymb}
\usepackage{mathtools}
\usepackage[shortlabels]{enumitem}
\newtheorem{lemma}{Lemma}
\newtheorem{theorem}{Theorem}
\newcommand{\F}{\mathbb F}
\newcommand{\GL}{\operatorname{GL}}
\newcommand{\circuittag}[1]{\text{[circuit $#1$]}}
\title{Four-Entropic Matroids Are Quaternary}
\author{Mohammad Hossein Kalantari, Shahram Khazaei}
\date{Source: arXiv:2608.20553v2 (CC BY 4.0)}
\begin{document}
\maketitle

\noindent\emph{Source and licence.} Four-Entropic Matroids Are Quaternary. Version arXiv:2608.20553v2 (CC BY 4.0), distributed by arXiv under the Creative Commons Attribution 4.0 International licence, http://creativecommons.org/licenses/by/4.0/, as recorded in the arXiv licence metadata for this exact version and shown on the abstract page https://arxiv.org/abs/2608.20553v2. Full licence notice: \texttt{licenses/arxiv-2608.20553-CC-BY-4.0.txt}.\par\smallskip

\noindent\textit{Editorial scope.} Counted unit: the article's theorem lem:p6 (source0.tex lines 655--657). The article's complete original proof of it is delivered with it (source0.tex lines 659--754, one \texttt{proof} environment of 96 lines). No mathematics is written, repaired or re-derived: the statement and the proof are the article's own lines, in the article's own notation, and no retained line is substituted or re-flowed.  Retained uncounted as the source-local context the proof needs: lines 541--558; lines 606--615; lines 610--614.  Nothing was omitted from any retained span, so the package carries no omission record.  No retained line contains a citation, so the wrapper carries no bibliography and no bibliography entry.  No retained line contains a figure environment, an image-inclusion command or drawing code, so no image is delivered and the package declares no asset dependency.  Editorial formatting only, in addition: an AMS \texttt{amsart} preamble that restates the article's own declarations and macros used by the retained lines, namely the environments \texttt{lemma}, \texttt{theorem} on separate counters, together with the standard packages those lines need.  The retained environments print in this document as Lemma 1, Theorem 1 in order of appearance, whereas the article prints the same items with the numbering of its own full document.  The complete native source of the article as distributed in the arXiv e-print for this exact version is bundled unchanged as \texttt{sources/208185/source0.tex}.
\begin{align}
\Gamma_{0134}^{4}(x,y,z)
 &=\Gamma_{0234}^{4}\bigl(x,\Gamma_{012}^{2}(x,y),z\bigr),
 &&\circuittag{0234}\label{eq:p6g1}\\
\Gamma_{0135}^{5}(x,y,z)
 &=\Gamma_{0235}^{5}\bigl(x,\Gamma_{012}^{2}(x,y),z\bigr),
 &&\circuittag{0235}\label{eq:p6g2}\\
\Gamma_{0135}^{5}(x,y,z)
 &=\Gamma_{0145}^{5}\bigl(x,y,\Gamma_{0134}^{4}(x,y,z)\bigr),
 &&\circuittag{0145}\label{eq:p6g3}\\
\Gamma_{0135}^{5}(x,y,z)
 &=\Gamma_{0345}^{5}\bigl(x,z,\Gamma_{0134}^{4}(x,y,z)\bigr),
 &&\circuittag{0345}\label{eq:p6g4}\\
\Gamma_{0135}^{5}(x,y,z)
 &=\Gamma_{0245}^{5}\bigl(x,\Gamma_{012}^{2}(x,y),
       \Gamma_{0134}^{4}(x,y,z)\bigr),
 &&\circuittag{0245}.\label{eq:p6g5}
\end{align}

\begin{lemma}\label{lem:p6-ternary}
If the $P_6$ quasigroup equations have a solution over $V$, then they have a
solution $(\Lambda_C)_{C\in\mathcal C}$ and matrices
$R,S,C,F\in\GL_2(\F_2)$ such that, for every $x,y,z\in V$,
\begin{align}
\Lambda_{012}^{2}(x,y)&=x+y,\label{eq:p6normal1}\\
\Lambda_{0134}^{4}(x,y,z)&=x+yR+zC,\label{eq:p6normal2}\\
\Lambda_{0135}^{5}(x,y,z)&=x+yS+zF.\label{eq:p6normal3}
\end{align}
\end{lemma}

\begin{align}
\Lambda_{012}^{2}(x,y)&=x+y,\\
\Lambda_{0134}^{4}(x,y,z)&=x+yR+zC,\\
\Lambda_{0135}^{5}(x,y,z)&=x+yS+zF.
\end{align}

\begin{theorem}\label{lem:p6}
The matroid $P_6$ has no partition representation of degree four.
\end{theorem}

\begin{proof}
Assume to the contrary that $P_6$ has a partition representation of degree
four. Then its quasigroup equations have a solution over $V$. By
Lemma~\ref{lem:p6-ternary}, we may take a solution
$(\Lambda_C)_{C\in\mathcal C}$ satisfying
\eqref{eq:p6normal1}--\eqref{eq:p6normal3}. Put
\[
Q=\Lambda_{012}^{2},\qquad
f=\Lambda_{0134}^{4},\qquad
g=\Lambda_{0135}^{5}.
\]
Rewriting \eqref{eq:p6g1}--\eqref{eq:p6g5} for this solution gives
\begin{align}
f(x,y,z)&=\Lambda_{0234}^{4}\bigl(x,Q(x,y),z\bigr),
&&\circuittag{0234}\label{eq:p6l1}\\
g(x,y,z)&=\Lambda_{0235}^{5}\bigl(x,Q(x,y),z\bigr),
&&\circuittag{0235}\label{eq:p6l2}\\
g(x,y,z)&=\Lambda_{0145}^{5}\bigl(x,y,f(x,y,z)\bigr),
&&\circuittag{0145}\label{eq:p6l3}\\
g(x,y,z)&=\Lambda_{0345}^{5}\bigl(x,z,f(x,y,z)\bigr),
&&\circuittag{0345}\label{eq:p6l4}\\
g(x,y,z)&=\Lambda_{0245}^{5}\bigl(x,Q(x,y),f(x,y,z)\bigr),
&&\circuittag{0245}.\label{eq:p6l5}
\end{align}

In \eqref{eq:p6l1}, put $u=x+y$, so $y=u+x$. Then
\[
\Lambda_{0234}^{4}(x,u,z)=x(I+R)+uR+zC.
\]
Because $\Lambda_{0234}^{4}$ is a ternary quasigroup, $I+R$ is invertible.
Equation \eqref{eq:p6l2} similarly gives that $I+S$ is invertible. Note that 
$\GL_2(\F_2) \cong S_3$ as groups. Since $I+R$ and $I+S$ are invertible, it follows
that $R,S \neq I$. Additionally, $R$ and $S$ cannot have order two, since otherwise it follows $(I+R)^2 = (I+S)^2 = 0$.
Thus each is one of the two order-three elements of $\GL_2(\F_2)$.

Next use \eqref{eq:p6l4}. Write $w=f(x,y,z)$ and solve
\[
y=(w+x+zC)R^{-1}.
\]
Substitution into $g$ gives
\[
\Lambda_{0345}^5(x,z,w)=x(I+R^{-1}S)+wR^{-1}S+z(CR^{-1}S+F).
\]
Since this is the affine expression for the ternary quasigroup
$\Lambda_{0345}^{5}(x,z,w)$, both
\[
I+R^{-1}S\qquad\text{and}\qquad CR^{-1}S+F
\]
are invertible. The first condition is equivalent to invertibility of $R+S$,
so $R\ne S$. The two order-three elements of $\GL_2(\F_2)$ are inverses;
after naming one of them $R$, we have $S=R^2$ and $I + S = R$. Therefore
$R^{-1}=R^2$ and $R^{-1}S=R$, and the second condition from the same circuit
becomes
\begin{equation}\label{eq:p6condCR}
F+CR\in\GL_2(\F_2).
\end{equation}

Equation \eqref{eq:p6l3} supplies another condition. Solving
$w=f(x,y,z)$ for $z$ gives
\[
\Lambda_{0145}^5(x,y,w)=x(I+C^{-1}F)+y(S+RC^{-1}F)+wC^{-1}F.
\]
The coefficient of $x$ must be invertible, equivalently
\begin{equation}\label{eq:p6condC}
F+C\in\GL_2(\F_2).
\end{equation}
The coefficient of $y$ reproduces \eqref{eq:p6condCR}.

Finally, in \eqref{eq:p6l5} put $u=x+y$ and again write $w=f(x,y,z)$.
Then
\[
z=\bigl(w+x(I+R)+uR\bigr)C^{-1}.
\]
The coefficient of $x$ in the affine expression for
$\Lambda_{0245}^{5}(x,u,w)$ is
\[
I+S+(I+R)C^{-1}F
=R+R^2C^{-1}F
=R^2\bigl(R^2+C^{-1}F\bigr).
\]
It must be invertible, which is equivalent to
\begin{equation}\label{eq:p6condCR2}
F+CR^2\in\GL_2(\F_2).
\end{equation}

Put $X=C^{-1}F$. Conditions \eqref{eq:p6condC}, \eqref{eq:p6condCR}, and
\eqref{eq:p6condCR2}, obtained respectively from the circuits $0145$, $0345$,
and $0245$, say that
\[
X+I,\qquad X+R,\qquad X+R^2
\]
are all invertible. The group $\GL_2(\F_2)\cong S_3$ consists of the identity,
three involutions, and the two order-three elements $R,R^2$. If $X+I$ is
invertible, then $X$ is neither the identity nor an involution. Hence
$X\in\{R,R^2\}$. One of $X+R$ and $X+R^2$ is then zero, a contradiction.
\end{proof}

\end{document}
